\documentclass[11pt]{article}
\usepackage[margin=1in]{geometry}
\usepackage{amsmath,amssymb,amsthm}
\usepackage{graphicx}
\usepackage{microtype}
\usepackage[colorlinks=true,linkcolor=blue,citecolor=blue,urlcolor=blue]{hyperref}

\newtheorem{theorem}{Theorem}
\newtheorem{lemma}{Lemma}
\newcommand{\C}{\mathbb C}
\newcommand{\E}{\mathcal E}
\newcommand{\R}{\mathcal R}
\newcommand{\A}{\mathcal A}

\title{A nonzero fixed-product preserver that destroys invertibility}
\author{Candidate counterexample for arXiv:2201.12912, Problem 3.5}
\date{11 August 2026}

\begin{document}
\maketitle

\begin{center}
\fbox{\parbox{0.91\linewidth}{\textbf{Status.} Candidate full
counterexample, likely valid, requiring human review. The example is a
separable commutative unital complex Banach algebra; the map is a bounded
linear involution; and the prescribed elements satisfy $c=d\ne0$ and
$c^2=c$.}}
\end{center}

\section{Source question}

Julius asks the following in Problem 3.5 of
\cite{Julius2023}, on PDF page 6:
\begin{quote}
If $\phi$ satisfies the hypotheses of Theorem 3.3, does $d\ne0$ imply that
$\phi$ preserves invertibility?
\end{quote}
Here the hypotheses are that $\A$ is a unital Banach algebra with dense
invertible group and that $\phi:\A\to\A$ is a bijective continuous linear map
satisfying
\[
        ab=c\quad\Longrightarrow\quad \phi(a)\phi(b)=d
\]
for fixed $c,d\in\A$.

\begin{figure}[ht]
\centering
\includegraphics[width=\linewidth]{figures/open_problem_crop.png}
\caption{The context and exact statement of Problem 3.5 in
\cite[p.~6]{Julius2023}. This image was rendered and cropped from the local
source paper PDF.}
\end{figure}

\section{Proof intuition}

The condition becomes very weak at a zero product in an integral domain:
$ab=0$ then forces one factor to be zero, so any linear map automatically
takes that pair to another zero product. A direct product supplies a separate
coordinate in which the prescribed product is nonzero. Thus the two desired
features can coexist: the first coordinate guarantees $d\ne0$, while the
domain coordinate leaves enough linear freedom to send a unit to a nonunit.

What remains is to produce a unital Banach integral domain with dense units.
We do this explicitly. A rapidly decreasing convolution weight makes a
nonunital weighted power-series domain radical. In its unitization, an element
is invertible exactly when its scalar coordinate is nonzero, so the units are
dense. Swapping the unit with the first power-series monomial is then a
bounded linear involution that does not preserve invertibility.

\section{A dense-unit Banach integral domain}

Put $\rho_n=e^{-n^2}$ for $n\ge1$, and let
\[
 \R=\left\{a=(a_n)_{n\ge1}:\
       \|a\|_\rho:=\sum_{n\ge1}|a_n|\rho_n<\infty\right\}.
\]
Give $\R$ the Cauchy convolution product
\[
 (a*b)_n=\sum_{j=1}^{n-1}a_jb_{n-j}.
\]
Let $z=(1,0,0,\ldots)$, so $z^n$ is the unit vector in coordinate $n$.

\begin{lemma}\label{lem:R}
$\R$ is a separable commutative radical Banach algebra and an integral
domain.
\end{lemma}

\begin{proof}
Since
\[
 \rho_{m+n}=e^{-(m+n)^2}\le e^{-m^2}e^{-n^2}=\rho_m\rho_n,
\]
Tonelli's theorem gives
\[
 \|a*b\|_\rho
 \le \sum_{m,n\ge1}|a_m||b_n|\rho_m\rho_n
 =\|a\|_\rho\|b\|_\rho.
\]
Completeness and separability follow from weighted $\ell_1$.

If $a,b\ne0$, let $m$ and $n$ be their least nonzero indices. The
coefficient of $a*b$ at $m+n$ is exactly $a_mb_n\ne0$. Hence $\R$ has no
zero divisors.

It remains to check radicality. The spectral-radius formula gives
\[
 r(z)=\lim_{k\to\infty}\|z^k\|_\rho^{1/k}
     =\lim_{k\to\infty}e^{-k}=0.
\]
By spectral mapping, every polynomial $p(z)$ with $p(0)=0$ has spectral
radius zero. Such polynomials (the finitely supported sequences) are dense
in $\R$. Since $\R$ is commutative, the spectral radius is subadditive; thus,
for any $a\in\R$ and any such polynomial $p$,
\[
        r(a)\le r(a-p)+r(p)\le\|a-p\|_\rho.
\]
Letting $p\to a$ yields $r(a)=0$. Every element is therefore
quasinilpotent, which says precisely that the commutative Banach algebra
$\R$ is radical. Standard spectral-radius facts used here may be found in
\cite{Dales2000}.
\end{proof}

Let $\A=\R^{\#}=\C1\oplus\R$ be the unitization, with norm
$\|\lambda1+r\|=|\lambda|+\|r\|_\rho$.

\begin{lemma}\label{lem:A}
$\A$ is a separable commutative unital Banach integral domain, and
\[
     \A^{-1}=\{\lambda1+r:\lambda\ne0\}.
\]
In particular, its invertible group is dense.
\end{lemma}

\begin{proof}
If $\lambda\ne0$, radicality of $\R$ implies that
$1+\lambda^{-1}r$ is invertible in the unitization, so $\lambda1+r$ is
invertible. If $\lambda=0$, the quotient homomorphism
$\A\to\A/\R\cong\C$ sends $r$ to zero, so $r$ cannot be invertible. This
proves the displayed characterization and density follows by perturbing the
scalar coordinate.

For domainhood, suppose $xy=0$. If either factor has nonzero scalar
coordinate, it is a unit and the other factor is zero. If both scalar
coordinates vanish, then $x,y\in\R$, and Lemma~\ref{lem:R} gives $x=0$ or
$y=0$.
\end{proof}

Let
\[
 \R_0=\{(a_n)\in\R:a_1=0\}.
\]
The first-coordinate functional is bounded, so
$\A=\C1\oplus\C z\oplus\R_0$ is a topological direct sum. Define
\begin{equation}\label{eq:L}
 L(\lambda1+\mu z+r_0)=\mu1+\lambda z+r_0.
\end{equation}
Then $L$ is complex-linear, $L^2=I$, and it is bounded: indeed,
\[
 \|L(\lambda1+\mu z+r_0)\|
 =|\mu|+e^{-1}|\lambda|+\|r_0\|_\rho
 \le e\bigl(|\lambda|+e^{-1}|\mu|+\|r_0\|_\rho\bigr).
\]
In particular, $L$ is a Banach-space automorphism. But $L(1)=z$ is not
invertible by Lemma~\ref{lem:A}.

\section{The counterexample}

\begin{theorem}\label{thm:counterexample}
There exist a separable commutative unital complex Banach algebra $\E$ with
dense invertible group, a bounded bijective complex-linear map
$\Phi:\E\to\E$, and a nonzero idempotent $c=d\in\E$ such that
\[
        xy=c\quad\Longrightarrow\quad\Phi(x)\Phi(y)=d,
\]
while $\Phi$ does not preserve invertibility. Consequently Problem 3.5 of
\cite{Julius2023} has a negative answer.
\end{theorem}

\begin{proof}
Take the coordinatewise direct product
\[
       \E=\C\times\A,
       \qquad \|(\alpha,a)\|=|\alpha|+\|a\|.
\]
By Lemma~\ref{lem:A}, $\E$ is separable and commutative, and its invertibles
\[
       \E^{-1}=\C^{\times}\times\A^{-1}
\]
are dense. Define
\[
       \Phi(\alpha,a)=(\alpha,L(a)),
       \qquad c=d=(1,0_\A).
\]
The map $\Phi$ is a bounded linear involution, hence is continuous and
bijective. Also $c=d$ is a nonzero idempotent.

Suppose
\[
       (\alpha,a)(\beta,b)=c.
\]
Coordinatewise multiplication gives $\alpha\beta=1$ and $ab=0$. Since
$\A$ is an integral domain, $a=0$ or $b=0$. In either case
$L(a)L(b)=0$, and therefore
\[
 \Phi(\alpha,a)\Phi(\beta,b)
 =\bigl(\alpha\beta,L(a)L(b)\bigr)
 =(1,0_\A)=d.
\]
Thus the fixed-product hypothesis holds for every factorization of $c$.

Finally, the identity of $\E$ is $u=(1,1_\A)$, whereas
\[
       \Phi(u)=(1,z).
\]
The first element is invertible, but the second is not because $z$ is not
invertible in $\A$. Hence $\Phi$ does not preserve invertibility even though
$d\ne0$.
\end{proof}

\section{Verification and scope}

The proof is exact and does not use a finite computation. The points most
worth checking are: (i) the spectral-radius passage proving radicality;
(ii) that radicality makes precisely the nonzero-scalar elements of the
unitization invertible; and (iii) that domainhood converts the condition
$ab=0$ into the disjunction $a=0$ or $b=0$. These are recorded in detail in
the accompanying \texttt{verification.md}.

This counterexample answers only Problem 3.5. It does not address Problem
3.4 concerning $\mathcal B(X)$.

\paragraph{Bounded novelty check.}
On 11 August 2026, the local run indexes were searched by arXiv id, title,
and the core terms. External exact/close searches combined the wording of
Problem 3.5 with ``fixed product preserving'', ``counterexample'', ``integral
domain'', ``radical Banach algebra'', and ``dense invertibles''. The audit
included the source and publisher pages and search-indexed portions of
Julius's 2026 survey chapter on fixed-product preservers. No prior explicit
answer or matching construction was found. The publisher lists seven citing
works, not all of whose full texts were individually inspected. Novelty
confidence is therefore moderate; a specialist citation review is
recommended. This limitation does not affect the internal validity of the
counterexample.

\begin{thebibliography}{9}
\bibitem{Julius2023}
H. Julius,
\emph{Fixed product preserving mappings on Banach algebras},
J. Math. Anal. Appl. \textbf{517} (2023), 126615;
\href{https://arxiv.org/abs/2201.12912}{arXiv:2201.12912}.

\bibitem{Dales2000}
H. G. Dales,
\emph{Banach Algebras and Automatic Continuity},
London Mathematical Society Monographs, New Series 24,
Clarendon Press, Oxford, 2000.
\end{thebibliography}

\end{document}
